\documentclass[11pt]{article}
\usepackage[a4paper,margin=21mm]{geometry}
\usepackage[no-math]{fontspec}
\setmainfont{Noto Serif}
\newfontfamily\CJKfont{Noto Serif CJK SC}
\XeTeXlinebreaklocale "zh"
\XeTeXlinebreakskip=0pt plus 1pt
\usepackage{amsmath,amssymb,amsthm,amscd,graphicx,color}
\usepackage{xurl}
\usepackage[unicode,hidelinks]{hyperref}
\allowdisplaybreaks
\setlength\emergencystretch{3em}
\numberwithin{equation}{section}
\usepackage{xspace}

\DeclareMathOperator{\SG}{S}
\usepackage{mathtools}
\usepackage{enumerate}

\usepackage{tikz}
\usepackage{gensymb}

\usepackage{bm,scalerel}
\usepackage[config, labelfont={bf},labelsep=period]{caption,subfig}
\captionsetup[subfigure]{subrefformat=simple,labelformat=simple,listofformat=subsimple}
\renewcommand\thesubfigure{(\alph{subfigure})}
\usepackage{mathrsfs}
\usetikzlibrary{matrix,arrows,calc}
\usetikzlibrary{decorations.markings}
\usetikzlibrary{patterns}
\usetikzlibrary{snakes}

\numberwithin{equation}{section}

\newtheorem{Theorem}{Theorem}[section]
\newtheorem*{Theorem*}{Theorem}
\newtheorem{Corollary}[Theorem]{Corollary}
\newtheorem{Lemma}[Theorem]{Lemma}
\newtheorem{Proposition}[Theorem]{Proposition}
 { \theoremstyle{definition}
\newtheorem{Definition}[Theorem]{Definition}
\newtheorem{Note}[Theorem]{Note}
\newtheorem{Example}[Theorem]{Example}
\newtheorem{Remark}[Theorem]{Remark} }

\def\A{\text{\rm vN}(G_{\s,\q}(\HH_\R))}
\def\C{{\mathbb C}}
\def\R{{\mathbb R}}
\def\N{{\mathbb N}}
\def\Z{{\mathbb Z}}
\def\D{{\rm D}}
\def\DD{D_q}
\def\d{{\rm d}}
\def\6{\, {\rm d}}
\def\i{{\rm i}}
\def\p{p}
\def\ri{{\rm i}}
\def\im{{\sf Im}}
\def\re{{\sf Re}}
\def\supp{{\sf supp}}
\def\B{b_{\s,\q}}
\def\G{G_{\s,\q}}
\def\GG{G^D_q}
\def\r{b}
\def\ell {n}
\def\F{\mathcal{F}_{\rm fin}(\HH)}
\def\id{I}
\def\state {\varphi}
\def\m{\mu_{\s,\q,x,y}}
\def\p{p_{\s,q,x}}
\def\qMP{{\rm MP}}
\def\P{\mathcal{P}}
\def\NC{\mathcal{NC}}
\def\Int{\mathcal{I}}
\def\Out{\rm Out}
\def\BL{{\mathbf{B}}}
\def\In{\rm In}
\def\PB{\mathcal{P}^{B}}
\def\PA{\mathcal{P}^{A}}
\def\Sign{\text{\normalfont sign}}
\def\Cr{\text{\normalfont Cr}}
\def\InS{\text{\normalfont Cs}}
\def\InNB{\text{\normalfont CNB}}
\def\NB{\text{\normalfont Nb}}
\def\SLNB{\text{\normalfont SLNB}}
\def\SL{\text{\normalfont SL}}
\def\Pair{\text{\normalfont Pair}}
\def\Sing{\text{\normalfont Sing}}
\renewcommand{\epsilon}{\varepsilon}
\def\Cov{\text{\normalfont Cover}}
\def\H{{\bf \mathcal K}_{ n}}
\newcommand{\e}{{\epsilon}}
\def\HH{{\bf \mathcal K}}
\def\HHE{{ \mathcal H}}
\def\inv{{pinv}}
\def\ninv{{ninv}}
\newcommand{\q}{q}
\newcommand{\s}{\alpha}%{q_{_{\text{-}}}}
\newcommand{\x}{\mathbf{x}_{\overline n} \otimes \mathbf{x}_n}
\DeclareMathOperator{\Part}{\mathcal{P}}
\DeclareMathOperator{\Span}{span}

\def\mylongmapsto#1{%
	\begin{tikzpicture}
		\draw (0,0.5mm) -- (0,-0.5mm);
		\newlength\mylength
		\setlength{\mylength}{\widthof{#1}}
		\draw[->] (0,0) -- (1.2\mylength,0) node[above,midway] {#1};
	\end{tikzpicture}
}

\newcommand\MatchingMeandersab[3]{%
	\begin{tikzpicture}[scale=0.4]
		% \draw(-0.5,0) -- ++ (#1+1,0);
		\foreach \x in {1,...,#1}{
			\draw[circle,fill] (\x,0)circle[radius=1mm]node[below]{};
		}
		\foreach \x/\y in {#2} {
			\pgfmathsetmacro{\Radius}{\y/2-\x/2}
			\draw(\x,0) arc[radius=\Radius, start angle=180, end angle=0];
			;
		}
		\foreach \x/\y in {#3} {
			\pgfmathsetmacro{\Radius}{\y/2-\x/2}
			\draw(\x,0) arc[radius=\Radius, start angle=-180, end angle=0];
			;}
		\foreach \x in {-#1,...,-1}{
			\draw[circle,fill] (\x,0)circle[radius=1mm]node[below]{};
		}
		\node at (-4,-0.7) { $\overline 4$};
		\node at (-3,-0.7) { $ \overline 3$};
		\node at (-2,-0.7) { $\overline 2$};
		\node at (-1,-0.7) { $\overline 1$};
		\node at (4,-0.7) { $ 4$};
		\node at (3,-0.7) { $ 3$};
		\node at (2,-0.7) { $ 2$};
		\node at (1,-0.7) { $ 1$};
	\end{tikzpicture}
}

\newcommand\MatchingMeandersabc[3]{%
	\begin{tikzpicture}[scale=0.7]
		\foreach \x in {1,...,#1}{
			\draw[circle,fill] (\x,0)circle[radius=1mm]node[below]{};
		}
		\foreach \x/\y in {#2} {
			\pgfmathsetmacro{\Radius}{\y/2-\x/2}
			\draw(\x,0) arc[radius=\Radius, start angle=180, end angle=0];
			;
		}
		\foreach \x/\y in {#3} {
			\pgfmathsetmacro{\Radius}{\y/2-\x/2}
			\draw(\x,0) arc[radius=\Radius, start angle=-180, end angle=0];
			;}
		\foreach \x in {-#1,...,-1}{
			\draw[circle,fill] (\x,0)circle[radius=1mm]node[below]{};
		}
		\node at (-4,-0.5) { $\overline 4$};
		\node at (-3,-0.5) { $ \overline 3$};
		\node at (-2,-0.5) { $\overline 2$};
		\node at (-1,-0.5) { $ \overline 1$};
		\node at (4,-0.5) { $ 4$};
		\node at (3,-0.5) { $ 3$};
		\node at (2,-0.5) { $ 2$};
		\node at (1,-0.5) { $ 1$};
	\end{tikzpicture}
}

\newcommand\MatchingMeandersexample[3]{%
	\begin{tikzpicture}[scale=0.2]
		\foreach \x in {1,...,#1}{
			\draw[circle,fill] (\x,0)circle[radius=1mm]node[below]{};
		}
		\foreach \x/\y in {#2} {
			\pgfmathsetmacro{\Radius}{\y/2-\x/2}
			\draw(\x,0) arc[radius=\Radius, start angle=180, end angle=0];
			;
		}
		\foreach \x/\y in {#3} {
			\pgfmathsetmacro{\Radius}{\y/2-\x/2}
			\draw(\x,0) arc[radius=\Radius, start angle=-180, end angle=0];
			;}
		\foreach \x in {-#1,...,-1}{
			\draw[circle,fill] (\x,0)circle[radius=1mm]node[below]{};
		}
	\end{tikzpicture}
}

\newcommand\MatchingMeandersPositiveTracea[3]{%
	\begin{tikzpicture}[scale=0.35]
		\foreach \x in {1,...,#1}{
			\draw[circle,fill] (\x,0)circle[radius=1mm]node[below]{};
		}
		\foreach \x/\y in {#2} {
			\pgfmathsetmacro{\Radius}{\y/2-\x/2}
			\draw(\x,0) arc[radius=\Radius, start angle=180, end angle=0];
			;
		}
		\foreach \x/\y in {#3} {
			\pgfmathsetmacro{\Radius}{\y/2-\x/2}
			\draw(\x,0) arc[radius=\Radius, start angle=-180, end angle=0];
			;}
		\foreach \x in {-#1,...,-1}{
			\draw[circle,fill] (\x,0)circle[radius=1mm]node[below]{};
		}
		\node at (-6.7,-0.6) { $\scriptscriptstyle{\bar v_m}$};
		\node at (-3.8,-0.6) { $\scriptscriptstyle{ \overline v_{i}}$};
		\node at (-0.8,-0.6) { $\scriptscriptstyle{\bar v_{1}}$};
		
		\node at (7.4,-0.6) { $\scriptscriptstyle{ v_m}$};
		\node at (4.2,-0.6) { $ \scriptscriptstyle{ v_{i}}$};
		\node at (1.2,-0.6) { $\scriptscriptstyle{ v_1}$};
		\node at (-2.5,4.5) { $\scriptscriptstyle{ \overline{W}_j}$};
		\node at (2.5,4.5) { $\scriptscriptstyle{ {W}_j}$};
		\node at (9,4) { $ \mapsto$};
	\end{tikzpicture}
}

\newcommand\MatchingMeandersPositiveTraceb[3]{%
	\begin{tikzpicture}[scale=0.35]
		\foreach \x in {1,...,#1}{
			\draw[circle,fill] (\x,0)circle[radius=1mm]node[below]{};
		}
		\foreach \x/\y in {#2} {
			\pgfmathsetmacro{\Radius}{\y/2-\x/2}
			\draw(\x,0) arc[radius=\Radius, start angle=180, end angle=0];
			;
		}
		\foreach \x/\y in {#3} {
			\pgfmathsetmacro{\Radius}{\y/2-\x/2}
			\draw(\x,0) arc[radius=\Radius, start angle=-180, end angle=0];
			;}
		\foreach \x in {-#1,...,-1}{
			\draw[circle,fill] (\x,0)circle[radius=1mm]node[below]{};
		}
		\node at (-6.7,-0.6) { $\scriptscriptstyle{\bar v_m}$};
		\node at (-3.8,-0.6) { $\scriptscriptstyle{ \overline v_{i}}$};
		\node at (-0.8,-0.6) { $\scriptscriptstyle{\bar v_{1}}$};
		
		\node at (7.4,-0.6) { $\scriptscriptstyle{ v_m}$};
		\node at (4.2,-0.6) { $ \scriptscriptstyle{ v_{i}}$};
		\node at (1.2,-0.6) { $\scriptscriptstyle{ v_1}$};
		\node at (-10,-0.6) { $\scriptscriptstyle{\overline{k+1}}$};
		\node at (10,-0.6) { $\scriptscriptstyle{k+1}$};
		\node at (-2.5,4.5) { $\scriptscriptstyle{ \overline{W}_j}$};
		\node at (2.5,4.5) { $\scriptscriptstyle{ {W}_j}$};
	\end{tikzpicture}
	
}

\newcommand\MatchingMeandersNegativeTracea[3]{%
	\begin{tikzpicture}[scale=0.5]
		\foreach \x in {1,...,#1}{
			\draw[circle,fill] (\x,0)circle[radius=1mm]node[below]{};
		}
		\foreach \x/\y in {#2} {
			\pgfmathsetmacro{\Radius}{\y/2-\x/2}
			\draw(\x,0) arc[radius=\Radius, start angle=180, end angle=0];
			;
		}
		\foreach \x/\y in {#3} {
			\pgfmathsetmacro{\Radius}{\y/2-\x/2}
			\draw(\x,0) arc[radius=\Radius, start angle=-180, end angle=0];
			;}
		\foreach \x in {-#1,...,-1}{
			\draw[circle,fill] (\x,0)circle[radius=1mm]node[below]{};
		}
		\node at (-6,-0.6) { $\overline a_1$};
		\node at (-3,-0.6) { $ \overline b_1$};
		\node at (3,-0.6) { $ b_1$};
		\node at (6,-0.6) { $a_1$};
		\node at (8,4) { $\pi \xrightarrow{\text{trace action }}\tilde{\pi}$};
	\end{tikzpicture}
}

\newcommand\MatchingMeandersNegativeTraceb[3]{%
	\begin{tikzpicture}[scale=0.5]
		\foreach \x in {1,...,#1}{
			\draw[circle,fill] (\x,0)circle[radius=1mm]node[below]{};
		}
		\foreach \x/\y in {#2} {
			\pgfmathsetmacro{\Radius}{\y/2-\x/2}
			\draw(\x,0) arc[radius=\Radius, start angle=180, end angle=0];
			;
		}
		\foreach \x/\y in {#3} {
			\pgfmathsetmacro{\Radius}{\y/2-\x/2}
			\draw(\x,0) arc[radius=\Radius, start angle=-180, end angle=0];
			;}
		\foreach \x in {-#1,...,-1}{
			\draw[circle,fill] (\x,0)circle[radius=1mm]node[below]{};
		}
		\node at (-1,-0.6) { $\overline a_1$};
		\node at (-4,-0.6) { $ \overline b_1$};
		\node at (4,-0.6) { $ b_1$};
		\node at (1,-0.6) { $a_1$};
	\end{tikzpicture}
	
}

\newcommand\MatchingMeanders[2]{%
	\begin{tikzpicture}[scale=0.7]
		\foreach \x in {1,...,#1}{
			\draw[circle,fill] (\x,0)circle[radius=1mm]node[below]{$ \x$};
		}
		\foreach \x/\y in {#2} {
			\pgfmathsetmacro{\Radius}{\y/2-\x/2}
			\draw(\x,0) arc[radius=\Radius, start angle=180, end angle=0];
			;
			\node at (0,5.45) { $\scriptstyle{\blacklozenge}$};
			\node at (0,4.2) { $\scriptstyle{\blacklozenge}$};
			\node at (0,1.7) { $\scriptstyle{\blacklozenge}$};
			\draw[] (0,0) -- (0,6);
		}
		
		\foreach \x in {#1,...,1}{
			\draw[circle,fill] (-\x,0)circle[radius=1mm]node[below]{$ \overline{\x}$};
		}
	\end{tikzpicture}
}

% Narrow formatting-only resolver: four byte-bound author diagram macros.
\makeatletter
\newcount\nativeOrphanActive
\newcount\nativeOrphanTotal
\newcount\nativeOrphanOne
\newcount\nativeOrphanTwo
\newcount\nativeOrphanThree
\newcount\nativeOrphanFour
\def\nativeCheckDiagramArgs#1#2#3#4{%
 \edef\nativeActual{\detokenize{#1|#2|#3}}%
%
 \edef\nativeExpected{\detokenize{#4}}%
 \ifx\nativeActual\nativeExpected\else\PackageError{nativecompat}{Unapproved author diagram arguments}{Only the exact source-bound calls are accepted}\fi}
\def\nativeRenderOrphan#1{\global\advance\nativeOrphanTotal by1\ifcase#1\PackageError{nativecompat}{Unknown render location}{Only source144,192,222,304}\or\global\advance\nativeOrphanOne by1\or\global\advance\nativeOrphanTwo by1\or\global\advance\nativeOrphanThree by1\or\global\advance\nativeOrphanFour by1\else\PackageError{nativecompat}{Unknown render location}{Only four macros}\fi\ifnum\nativeOrphanTotal>32\PackageError{nativecompat}{Too many orphan delimiters}{Expected32}\fi}

\renewcommand\MatchingMeandersabc[3]{%
	\begin{tikzpicture}[scale=0.7]
		\foreach \x in {1,...,#1}{
			\draw[circle,fill] (\x,0)circle[radius=1mm]node[below]{};
		}
		\foreach \x/\y in {#2} {
			\pgfmathsetmacro{\Radius}{\y/2-\x/2}
			\draw(\x,0) arc[radius=\Radius, start angle=180, end angle=0];
			\nativeRenderOrphan{1}
		}
		\foreach \x/\y in {#3} {
			\pgfmathsetmacro{\Radius}{\y/2-\x/2}
			\draw(\x,0) arc[radius=\Radius, start angle=-180, end angle=0];
			;}
		\foreach \x in {-#1,...,-1}{
			\draw[circle,fill] (\x,0)circle[radius=1mm]node[below]{};
		}
		\node at (-4,-0.5) { $\overline 4$};
		\node at (-3,-0.5) { $ \overline 3$};
		\node at (-2,-0.5) { $\overline 2$};
		\node at (-1,-0.5) { $ \overline 1$};
		\node at (4,-0.5) { $ 4$};
		\node at (3,-0.5) { $ 3$};
		\node at (2,-0.5) { $ 2$};
		\node at (1,-0.5) { $ 1$};
	\end{tikzpicture}
}
\renewcommand\MatchingMeanders[2]{%
	\begin{tikzpicture}[scale=0.7]
		\foreach \x in {1,...,#1}{
			\draw[circle,fill] (\x,0)circle[radius=1mm]node[below]{$ \x$};
		}
		\foreach \x/\y in {#2} {
			\pgfmathsetmacro{\Radius}{\y/2-\x/2}
			\draw(\x,0) arc[radius=\Radius, start angle=180, end angle=0];
			\nativeRenderOrphan{2}
			\node at (0,5.45) { $\scriptstyle{\blacklozenge}$};
			\node at (0,4.2) { $\scriptstyle{\blacklozenge}$};
			\node at (0,1.7) { $\scriptstyle{\blacklozenge}$};
			\draw[] (0,0) -- (0,6);
		}
		
		\foreach \x in {#1,...,1}{
			\draw[circle,fill] (-\x,0)circle[radius=1mm]node[below]{$ \overline{\x}$};
		}
	\end{tikzpicture}
}
\renewcommand\MatchingMeandersPositiveTracea[3]{%
	\begin{tikzpicture}[scale=0.35]
		\foreach \x in {1,...,#1}{
			\draw[circle,fill] (\x,0)circle[radius=1mm]node[below]{};
		}
		\foreach \x/\y in {#2} {
			\pgfmathsetmacro{\Radius}{\y/2-\x/2}
			\draw(\x,0) arc[radius=\Radius, start angle=180, end angle=0];
			\nativeRenderOrphan{3}
		}
		\foreach \x/\y in {#3} {
			\pgfmathsetmacro{\Radius}{\y/2-\x/2}
			\draw(\x,0) arc[radius=\Radius, start angle=-180, end angle=0];
			;}
		\foreach \x in {-#1,...,-1}{
			\draw[circle,fill] (\x,0)circle[radius=1mm]node[below]{};
		}
		\node at (-6.7,-0.6) { $\scriptscriptstyle{\bar v_m}$};
		\node at (-3.8,-0.6) { $\scriptscriptstyle{ \overline v_{i}}$};
		\node at (-0.8,-0.6) { $\scriptscriptstyle{\bar v_{1}}$};
		
		\node at (7.4,-0.6) { $\scriptscriptstyle{ v_m}$};
		\node at (4.2,-0.6) { $ \scriptscriptstyle{ v_{i}}$};
		\node at (1.2,-0.6) { $\scriptscriptstyle{ v_1}$};
		\node at (-2.5,4.5) { $\scriptscriptstyle{ \overline{W}_j}$};
		\node at (2.5,4.5) { $\scriptscriptstyle{ {W}_j}$};
		\node at (9,4) { $ \mapsto$};
	\end{tikzpicture}
}
\renewcommand\MatchingMeandersPositiveTraceb[3]{%
	\begin{tikzpicture}[scale=0.35]
		\foreach \x in {1,...,#1}{
			\draw[circle,fill] (\x,0)circle[radius=1mm]node[below]{};
		}
		\foreach \x/\y in {#2} {
			\pgfmathsetmacro{\Radius}{\y/2-\x/2}
			\draw(\x,0) arc[radius=\Radius, start angle=180, end angle=0];
			\nativeRenderOrphan{4}
		}
		\foreach \x/\y in {#3} {
			\pgfmathsetmacro{\Radius}{\y/2-\x/2}
			\draw(\x,0) arc[radius=\Radius, start angle=-180, end angle=0];
			;}
		\foreach \x in {-#1,...,-1}{
			\draw[circle,fill] (\x,0)circle[radius=1mm]node[below]{};
		}
		\node at (-6.7,-0.6) { $\scriptscriptstyle{\bar v_m}$};
		\node at (-3.8,-0.6) { $\scriptscriptstyle{ \overline v_{i}}$};
		\node at (-0.8,-0.6) { $\scriptscriptstyle{\bar v_{1}}$};
		
		\node at (7.4,-0.6) { $\scriptscriptstyle{ v_m}$};
		\node at (4.2,-0.6) { $ \scriptscriptstyle{ v_{i}}$};
		\node at (1.2,-0.6) { $\scriptscriptstyle{ v_1}$};
		\node at (-10,-0.6) { $\scriptscriptstyle{\overline{k+1}}$};
		\node at (10,-0.6) { $\scriptscriptstyle{k+1}$};
		\node at (-2.5,4.5) { $\scriptscriptstyle{ \overline{W}_j}$};
		\node at (2.5,4.5) { $\scriptscriptstyle{ {W}_j}$};
	\end{tikzpicture}
	
}
\let\nativeOrigAbc\MatchingMeandersabc
\renewcommand\MatchingMeandersabc[3]{\begingroup\nativeOrphanActive=1\nativeCheckDiagramArgs{#1}{#2}{#3}{4|-4/1, -3/-2,2/3,-1/4 |}\nativeOrigAbc{#1}{#2}{#3}\endgroup}
\let\nativeOrigMeanders\MatchingMeanders
\renewcommand\MatchingMeanders[2]{\begingroup\nativeOrphanActive=2\nativeCheckDiagramArgs{#1}{#2}{}{10|-1/3, -3/1 ,-6/5,-5/6,7/10,-10/-7,-2/9,-9/2|}\nativeOrigMeanders{#1}{#2}\endgroup}
\let\nativeOrigTraceA\MatchingMeandersPositiveTracea
\renewcommand\MatchingMeandersPositiveTracea[3]{\begingroup\nativeOrphanActive=3\nativeCheckDiagramArgs{#1}{#2}{#3}{9|-9/-8,8/9,-6/2,-2/6 |}\nativeOrigTraceA{#1}{#2}{#3}\endgroup}
\let\nativeOrigTraceB\MatchingMeandersPositiveTraceb
\renewcommand\MatchingMeandersPositiveTraceb[3]{\begingroup\nativeOrphanActive=4%
 \edef\nativeActual{\detokenize{#1|#2|#3}}%
 \def\nativeExpected{10|-9/-8,8/9,-6/2,-2/6 ,-10/-4,4/10 |}%
 \ifx\nativeActual\nativeExpected\else\def\nativeExpected{10|-9/-8,8/9,-6/2,-2/6 ,-10/4,-4/10 |}%
 \ifx\nativeActual\nativeExpected\else\PackageError{nativecompat}{Unapproved trace-B arguments}{Only the two exact original positive and negative paths are accepted}\fi\fi%
 \nativeOrigTraceB{#1}{#2}{#3}\endgroup}
\AtEndDocument{%
 \ifnum\nativeOrphanTotal=32\else\PackageError{nativecompat}{Orphan total mismatch}{Expected32}\fi%
 \ifnum\nativeOrphanOne=4\else\PackageError{nativecompat}{Abc count mismatch}{Expected4}\fi%
 \ifnum\nativeOrphanTwo=8\else\PackageError{nativecompat}{Meanders count mismatch}{Expected8}\fi%
 \ifnum\nativeOrphanThree=8\else\PackageError{nativecompat}{Trace-A count mismatch}{Expected8}\fi%
 \ifnum\nativeOrphanFour=12\else\PackageError{nativecompat}{Trace-B count mismatch}{Expected12}\fi%
 \typeout{NATIVE-COMPAT: four exact diagrams consumed32 orphan delimiters; counts4,8,8,12}}
\makeatother

\makeatletter\newlabel{figexample1}{{6}{}{Original source locator}{}{}}\makeatother
\begin{document}
\begin{center}\Large Type-B double-Fock Wick formula with crossing and negative-pair weights for real simple tensors\end{center}
\noindent\textbf{Stable ID:} 1146\quad\textbf{Author:} Marek Bozejko; Wiktor Ejsmont\par
\noindent\textbf{Work:} The Double Fock Space of Type B\par
\noindent\textbf{Source:} \url{https://sigma-journal.com/2023/040/}\par
\noindent\textbf{Version:} Official SIGMA 19 (2023) article 040, DOI 10.3842/SIGMA.2023.040; journal PDF and native TeX acquired 2026-09-30, exact bytes pinned by SHA256\par
\noindent\textbf{Rights:} CC BY 4.0. Authors retain copyright. \url{https://creativecommons.org/licenses/by/4.0/}. Reuse and typesetting adaptation; original language and mathematical content preserved.\par
\noindent\textbf{Difficulty (prior selection preserved):} {\CJKfont H1: The proof constructs a weighted correspondence between double creation/annihilation histories and reflection-symmetric pair-and-singleton partitions. In the negative and positive annihilation cases it tracks crossings and covered singletons separately, obtaining the different powers of q and the negative-pair power of alpha. A bijective induction then removes the singletons under vacuum expectation to establish the full joint moment formula, rather than calculating a single Gaussian marginal.}\par
\medskip\noindent\textbf{Editorial boundary note (not author text):} Complete original Theorem4.6 and its full creation/positive/negative annihilation correspondence, weighted bijective induction and exhaustive final vacuum pairing sum. All real-Hilbert/complexification/simple-tensor and alpha,q in(-1,1) hypotheses, Coxeter action/length definitions, double symmetrization and free/adjoint operator conventions, exact parabolic decomposition and explicit positive/negative annihilation formulas, Gaussian/vacuum state definitions, reflection-symmetric pairs/singletons and every crossing/covered-singleton/negative-pair statistic are retained. Prior Coxeter positivity and parabolic decomposition are exact author-stated formal prerequisites; the current adjoint formula is retained with the exact cited elementary method, without editor derivation. The omitted norm estimates and one-variable orthogonal-polynomial theorem are not inputs to this algebraic vacuum identity. Every essential native TikZ induction diagram and label remains editable. One primary multidimensional moment outcome expressly motivated in the introduction, with supporting construction/statistics receiving no extra counts. Unused bbm import is omitted; no bbm command occurs in the source. The unused stackengine import is omitted after command inventory found no stackengine command; scalerel remains available from its original bm/scalerel import. Root-approved render-only resolver consumes exactly32 orphan second arc delimiters only inside the four source-bound MatchingMeandersabc/MatchingMeanders/PositiveTracea/PositiveTraceb macros, rejects unapproved arguments/locations and checks per-macro4/8/8/12 totals. Literal source paths and labels are unchanged. Named Noto Serif prose font preserves the original literal bibliography cedilla; standard math fonts are unchanged.\par
\noindent\textbf{External original reference figexample1:} Original fourth-moment example Figure6, source1099–1142. Mentioned only in the positive-pair counting remark, not an omitted selected proof diagram. The selected proof retains all its native positive/negative induction diagrams.\par
\medskip\hrule\medskip\noindent\textbf{Author text begins below.}\par
\setcounter{section}{1}
% Original source lines 440--490
\section{Preliminaries }
\label{sec:preliminaries}
In the present section we show that the construction of type B Fock space \cite{BozejkoEjsmontHasebe2015} can be adapted to recover the
double Fock space of type B.
In the following we will briefly describe the tools we use in this investigation,
which is partially contained in the previous paper \cite{BozejkoEjsmontHasebe2015}.
For further information the reader is referred to \cite{BozejkoEjsmontHasebe2015} and the references therein.

\subsection{Coxeter groups of type B}
Recall that the Coxeter
group of type B (also known as the hyperoctahedral group) of degree $n$, denoted by $B(n)$, is generated
by the elements $\pi_0, \pi_1, \dots , \pi_{n-1}$ subject to the defining relations $\pi_i^2=e$, $0\leq i \leq n-1$, $(\pi_0\pi_1)^4=(\pi_i \pi_{i+1})^3=e$, $1\leq i < n-1$ and $(\pi_i \pi_j)^2=e$ if $|i-j|\geq2$, $0\leq i,j\leq n-1$. Note that $\{\pi_i\mid i=1,\dots,n-1\}$ generate the symmetric group $S(n)$.
The Coxeter diagram for $B(n)$ is described in Figure \ref{fig:BN}.
\begin{figure}[htp]
\centering
		\begin{tikzpicture}
			[scale=.5,auto=left,every node/.style={circle}]
			\node (n7) at (1,3.6) {$\pi_{n-1}$};
			\node (n6) at (1,3) {$\circ$};
			\node (n5) at (-3,3) {$\dots$};
			\node (n3) at (-7,3.6) {$\pi_2$};
			\node (n1) at (-7,3) {$\circ$};
			\node (n4) at (-11,3.6) {$\pi_{1}$};
			\node (n2) at (-11,3) {$\circ$};
			\node (n8) at (-15,3.6) {$\pi_{0}$};
			\node (n0) at (-15,3) {$\circ$};
			
			\foreach \from/\to in {n2/n1, n1/n2,n1/n5,n5/n6}
			\draw (\from) -- (\to);
			\draw (-14.5,3.1) -- (-11.5,3.1);
			\draw (-14.5,2.9) -- (-11.5,2.9);
		\end{tikzpicture}
		\caption{The Coxeter diagram for $B(n)$.}
		\label{fig:BN}
\end{figure}


We express $\sigma \in B(n)$ in an irreducible form
\[
\sigma=\pi_{i_1} \cdots \pi_{i_{k}}, \qquad 0\leq i_1,\dots,i_k \leq n-1,
\]
i.e., in a form with the minimal length, and in this case let
\begin{align*}
	&l_1(\sigma)= \text{the number of the occurrences of the factor $\pi_0$ in $\sigma$},
\\[1mm]
	&l_2(\sigma) = \text{the number of the occurrences of all the factor of the form $\pi_1,\dots,\pi_{n-1}$ in $\sigma$}.
\end{align*}

\begin{Remark}
These definitions do not depend on the way we express $\sigma$ in an irreducible form, and therefore, $l_1(\sigma)$ and $l_2(\sigma)$ are well defined (see \cite[Proposition~1]{BozejkoSzwarc2003}).
\end{Remark}

\setcounter{section}{2}
% Original source lines 567--722
\section{The double Fock space of type B}
%\label{subsec:Fock}

Let $H_\R$ be a separable real Hilbert space and let $H$ be its
complexification with the inner product~$\langle\cdot,\cdot\rangle$, linear in the right component and anti-linear in the left.
The Hilbert space $\HH:=H\otimes {H}$ is the complexification of its
real subspace $\HH_\R:=H_\R\otimes {H}_\R$, with the inner product
\[
\langle x\otimes y,\xi\otimes \eta \rangle_{\HH} = \langle
x,\xi\rangle\langle y,\eta\rangle.
\]

We define $\H:=H^{\otimes n} \otimes {{H}^{\otimes n}}=H^{\otimes 2n}$ and instead of
indexing its simple tensors by $\{1,\dots,2n\}$ we will index them by $[\pm n]=\{-n,\dots,-1,1,\dots,n\}$ and we will use the
typical involution $\bar n=-n$ for $n\in \N$
\[
\H\ni\mathbf{x}_{\overline n} \otimes \mathbf{x}_n=x_{\overline n}\otimes\cdots \otimes x_{\overline 1} \otimes x_1\otimes\cdots \otimes x_n=x_{\overline n}\otimes \cdots \otimes x_{n}.
\]
We use this convention for indexing the elements of $\H$ to define a natural action of the
hyperoctahedral group $B(n)$ on $\H$ by setting
\begin{align*}
\sigma\colon\, &\H\to \H,
\\
	&x_{\overline n}\otimes\cdots \otimes x_{\overline 1} \otimes x_1\otimes\cdots \otimes x_n\mapsto x_{\sigma^{-1}(\overline n)}\otimes\cdots \otimes x_{\sigma^{-1}(\overline 1)} \otimes x_{\sigma^{-1}(1)}\otimes\cdots \otimes x_{\sigma^{-1}(n)}
\end{align*}
for any $\sigma \in B(n)$.

\begin{Remark}
	It is worth to mention that in the present paper we act on the indexes, but in the previous one, \cite{BozejkoEjsmontHasebe2015}, we were acting on the vectors.
\end{Remark}



Let $\F$ be the algebraic full Fock space over $\HH$:
\begin{equation*}
	\F:= \bigoplus_{n=0}^\infty \H= \bigoplus_{n=0}^\infty H^{\otimes 2n},
\end{equation*}
with the convention that $\HH^{\otimes 0} =H^{\otimes 0} \otimes H^{\otimes 0}=\C\Omega \otimes \Omega$ is the one-dimensional normed space along with the unit vector $\Omega \otimes \Omega$.
We equip $\F$ with the inner product
\[
\langle x_{\overline n} \otimes \cdots \otimes x_{\overline 1} \otimes x_1 \otimes \cdots \otimes x_n, y_{\overline m} \otimes \cdots \otimes y_{\overline 1} \otimes y_1 \otimes \cdots \otimes y_m\rangle_{0,0}:= \delta_{m,n}\prod_{i=\overline n }^n \langle x_i, y_i\rangle.
\]
By definition, the action of the generators $\pi_i$ for $0\leq i \leq n-1$ on
\begin{align*}
\eta = x_{\bar n} \otimes \cdots \otimes x_{\bar 1} \otimes x_1 \otimes \cdots \otimes x_{n}\in \H
\end{align*}
is given for $i\geq 1$ and $n\geq 2$ by
\begin{align*}
\pi_i(\eta) =
	x_{\bar n} \otimes \cdots \otimes x_{\overline{i}} \otimes x_{\overline{i+1}} \otimes \cdots \otimes x_{\bar 1} \otimes x_1 \otimes \cdots \otimes x_{i+1} \otimes x_{i} \otimes \cdots \otimes x_{n}, \end{align*}
and for $i=0$ and $n\geq 1$ by
\begin{align*}
\pi_0(\eta) = x_{\bar n} \otimes \cdots \otimes x_1 \otimes x_{\bar 1}\otimes \cdots \otimes x_{n}.
\end{align*}

For the parameters $ -1\leq \alpha, q \leq 1$, we define the symmetrization operators
\begin{align*}
	&P_{\s,\q}^{(n)}:= \sum_{\sigma \in B(n)}\s^{l_1(\sigma)} \q^{l_2(\sigma)} \, \sigma,\qquad n \geq1, \\
	&P_{\s,\q}^{(0)}:= \id_{H^{\otimes 0}\otimes H^{\otimes 0}}.
	\intertext{Moreover, let }
	&P_{\s,\q}:=\bigoplus_{n=0}^\infty P_{\s,\q}^{(n)}
\end{align*} be the \emph{type B symmetrization operator} acting on the algebraic full Fock space.
\begin{Remark}
	From
	Bo\.zejko and Speicher \cite[Theorem 2.1]{BozejkoSpeicher1994}, the operator $P_{\s,\q}^{(n)}$ is positive for $ -1\leq \alpha, q \leq 1$. If $ -1< \alpha, q <1$, then $P_{\s,\q}^{(n)}$ is a strictly positive operator meaning that it is positive and
	$\ker P_{\s,\q}^{(n)}=\{0\}$.
\end{Remark}
For $\x\in \HH_{ n}$ and $\mathbf{y}_{\overline m} \otimes \mathbf{y}_m\in
\HH_{ m}$, we deform the inner product $\langle\cdot,\cdot\rangle_{0,0}$ by using the type~B symmetrization operator:
\begin{align*}
	\langle \x,\mathbf{y}_{\overline m} \otimes \mathbf{y}_m\rangle_{\s,\q}:=\delta_{n,m}\langle \x,P_{\s,\q}^{(m)}\mathbf{y}_{\overline m} \otimes \mathbf{y}_m\rangle_{0,0}.
\end{align*}

For $x\in H$, let $l(x)$ and $r(x)$ be the free left and right
annihilation operators on $H^{\otimes n} $, respectively, defined by
the equations
\begin{gather*}
l^\ast(x)(x_1\otimes \dots \otimes x_n ):=x\otimes x_1\otimes \dots \otimes x_n ,
\\
l(x)(x_1\otimes \dots \otimes x_n ):=\langle x, x_1 \rangle x_2\otimes \dots \otimes x_n ,
\\
r^\ast(x)(x_1\otimes \dots \otimes x_n ):= x_1\otimes \dots \otimes x_n \otimes x,
\\
r(x)(x_1\otimes \dots \otimes x_n ):=\langle x, x_n \rangle x_1\otimes \dots \otimes x_{n-1},
\end{gather*}
where the adjoint is taken with respect to the free inner product.
The left-right creation and annihilation operators $\r^\ast(x\otimes y)$, $\r(x\otimes y)$ on $\F$ are defined as follows:
\begin{alignat*}{2}
&\r^\ast(x\otimes y)(\x):=l^\ast(x)\mathbf{x}_{\overline n}\otimes
	r^\ast( y) \mathbf{x}_n, \qquad &&\r^\ast(x\otimes y)\Omega \otimes \Omega:=x \otimes y,
\\
&\r(x\otimes y)(\x ):= l( x) \mathbf{x}_{\overline n}\otimes r(y)\mathbf{x}_n , \qquad &&\r(x\otimes y)\Omega \otimes \Omega :=0,
\end{alignat*}
where $\x\in \H$, $n\geq 1$.
It holds that $[\r^\ast(x\otimes y)]^\ast = \r(x\otimes y)$ where the adjoint is taken with respect to $\langle \cdot, \cdot \rangle_{0,0}$.

\begin{Definition} Let $\q,\s \in (-1,1)$. The algebraic full Fock space $\F$ equipped with the inner product $\langle\cdot,\cdot \rangle_{\s,\q}$ is called the \emph{double Fock space of type B} and denoted by $\mathcal{F}_{\s,\q}(\HH)$.
	For $x\otimes y \in \HH$ we define $\B^\ast(x\otimes y):=
	\r^\ast(x\otimes y)$ and we consider its adjoint operator $\B(x\otimes
	y)$ with
	respect to the inner product $\langle\cdot,\cdot \rangle_{\s,\q}$
	acting on the Hilbert space $\mathcal{F}_{\s,\q}(\HH)$. The operators $\B^\ast(x\otimes y)$ and $\B(x\otimes y)$ are called \emph{double creation and double annihilation operator of type B, respectively}.
\end{Definition}

\begin{Remark}
We would like to emphasize that our double Fock space of type B is different from~\cite{BozejkoEjsmontHasebe2015}.
	In the present version of the model, we apply a more natural operation on tensor product, which finally contributes to the more natural combinatorics of calculating the Gaussian moments.
	In~\cite{BozejkoEjsmontHasebe2015}, we put the action of the symmetrization on $n$ points but in the current ap\-proach we give $2n$ points which is close to our construction from the articles~\cite{BDEG2021,Ejsmont2020}. In~\cite{BozejkoEjsmontHasebe2015}, we defined the creator in the same way as in the case of type A, and then we did not obtain a~natural combinatorics which is compatible with partitions of the set $\{\pm 1,\dots, \pm n\}$.
	Now if we start our construction with the creator operator of double type B, the situation is completely different from~\cite{BozejkoEjsmontHasebe2015}.
\end{Remark}
The following proposition can be derived directly
from \cite[Proposition 2.3]{BozejkoEjsmontHasebe2015}.

\begin{Proposition}%\label{prop1}
We have the decomposition
\begin{equation*}%\label{decomposition}
P^{(n)}_{\s,\q}=\big( \id \otimes P^{(n-1)}_{\s,\q}\otimes \id \big)R^{(n)}_{\s,\q}\qquad \text{on}\quad \H, \qquad n\geq 1,
\end{equation*}
where\vspace{-2mm}
\begin{gather*}
R^{(n)}_{\s,\q} = \id+\sum_{k=1}^{n-1}\q^{k}\pi_{n-1}\cdots \pi_{n-k} + \s \q^{n-1}\pi_{n-1} \pi_{n-2} \cdots \pi_{1}\pi_0\bigg(1+\sum_{k=1}^{n-1}\q^{k}\pi_{1}\cdots \pi_{k}\bigg), \quad n\geq 2,
\\
R^{(1)}_{\s,\q} =\id+\s\pi_0.% \quad n=1.
\end{gather*}
\end{Proposition}

We can compute the annihilation operator in terms of $R_{\s,\q}^{(n)}$ by using the similar
method to that in \cite[Proposition~2.4]{BozejkoEjsmontHasebe2015}.

\begin{Proposition}%\label{prop2}
For $n \geq1$ and $x\otimes y\in \HH_\R$, we have\vspace{-1mm}
\begin{equation*}
\B(x\otimes y)=\r(x\otimes y) R^{(n)}_{\s,\q}\qquad \text{on}\quad \H.
\end{equation*}
\end{Proposition}

\begin{Remark}\label{thm1}
	Using the above notation, we can decompose $\B(x\otimes y)$ into the positive part $p_\q$ and the negative part $\ell_{\q}$ as\vspace{-1mm}
\[
	\B(x\otimes y)= p_\q(x\otimes y)+ \s \ell_{\q}(x\otimes y), \qquad x\otimes y \in \HH,
\]
where\vspace{-2mm}
\begin{align}\label{rq}
&p_\q(x\otimes y)\eta = \sum_{k=1}^n \q^{n-k}\langle x, x_{\bar{k}}\rangle \langle y, x_k \rangle\, x_{\bar n }\otimes \cdots \otimes \check{x}_{\bar k} \otimes \cdots \otimes x_{\bar 1} \nonumber
\\[-1mm]
&\hphantom{p_\q(x\otimes y)\eta = \sum_{k=1}^n}
{}\otimes x_1\otimes \cdots \otimes \check{x}_k \otimes \cdots \otimes x_n,
\\
&\ell_{\q}(x\otimes y)\eta =\q^{n-1}\sum_{k=1}^n \q^{k-1}\langle x, x_{{k}}\rangle \langle y, x_{\bar{k}}\rangle\, x_{\bar n }\otimes \cdots \otimes \check{x}_{\bar k} \otimes \cdots \otimes x_{\bar 1} \nonumber
\\
&\hphantom{\ell_{\q}(x\otimes y)\eta =\q^{n-1}\sum_{k=1}^n}
\otimes x_1\otimes \cdots \otimes \check{x}_k \otimes \cdots \otimes x_n,\label{lq}
	\end{align}
	where $\eta=x_{\bar n }\otimes \cdots \otimes x_n$.
	We called the operator $p_\q$ the positive part and $\ell_{\q}$ the negative part because in the next section we see that they contribute to positive and negative partitions of type B, respectively.
\end{Remark}


% Original source lines 769--782
\section{The double Gaussian operator of type B }
In this non-commutative setting, random variables are understood to be the elements of the
$\ast$-algebra generated by $\{\B(x\otimes y) ,\B^\ast(x\otimes y)\mid x\otimes y \in \HH_\R\}$. Particularly interesting are their mixed moments.
In order to work effectively on this object we need to define the corresponding statistics. We provide an explicit formula for the combinatorial moments, involving the number of crossings and negative pairs of a partition. First, we need to define the operators, the set of partitions and statistics of type B.

\subsection{The orthogonal polynomials}
\begin{Definition}
	The operator
	\begin{equation*}
		\G(x\otimes y)= \B(x\otimes y) +\B^\ast(x\otimes y),\qquad x\otimes y \in \HH_\R
	\end{equation*}
	on $\F$ is called the \emph{double Gaussian operator of type B}.
Denote by $\state$ the vacuum vector state $\state(\cdot)=\langle\Omega \otimes \Omega, \cdot\text{ } \Omega\otimes \Omega\rangle$.
\end{Definition}

\setcounter{subsection}{1}\setcounter{Theorem}{2}
% Original source lines 829--1097
\subsection{Pair partitions of type B}
%\label{subsec:PartitionsB}

Let $S$ be an ordered set. Then $\pi = \{ V_1,\dots, V_p\}$ is a partition of $S$, if the $V_i \neq \varnothing$ are
ordered and disjoint sets $V_i=(v_1,\dots,v_k)$, where $v_1<\dots<v_k$, whose union is $S$.
For $V \in \pi$ , we say that~$V$ is a \emph{block of $\pi$}.
Any partition $\pi$ defines an equivalence relation on $S$,
denoted by $\sim_\pi$, such that the equivalence classes are the
blocks of~$\pi$.
Therefore, $i\sim_\pi j$ if $i$ and $j$ belong to the same block of $\pi$.
A block of $\pi$ is called a \emph{singleton} if it consists of one element.
Similarly, a~block of $\pi$ is called a \emph{pair} if it consists of two elements.
Let $\Sing(\pi)$ and $\Pair(\pi)$ denote the set of all singletons and pairs of $\pi$, respectively.
The set of partitions of $S$ is denoted by $\Part(S)$, in the case where $S =
[n] := \{1, \dots , n\}$ we write $\Part(n)$ := $\Part([n])$.
We denote by $\P_{1,2}(S)$ the subset of partitions $\pi \in \Part(S)$ whose every block is either a
pair or a singleton.

From now on, we will work on a set $[\pm n]:=\{\bar n, \dots , \bar 1, 1,\dots, n\}$.
For a pair $V=(a,b)$ (or a~singleton $V=(a)$), we denote its reflection by $\overline V=(\bar b , \bar a)$ ($\overline V=( \bar a)$), where $\bar a =-a$.
Similarly, we define
\[
\bar \pi: =\big\{\overline{V}\mid V\in \pi,\text{ } \pi \in \P_{1,2}([\pm n])\big\}.
\]

\begin{Definition} \label{def:partycji}
	We denote by $\P_{1,2}^{B}(n)$ the subset of partitions $\pi \in \P_{1,2}([\pm n])$
	such that they are symmetric $\overline{\pi} =
	\pi$ (which is invariant under the bar operation), but every pair $V \in \pi$ is different from its reflection
	$\overline{V}$, i.e.,~$V \neq \overline{V}$.
	We call $\P_{1,2}^{B}(n)$ the set of partitions of type B.
\end{Definition}

From Definition~\ref{def:partycji} it follows that for every block in $B\in \pi$, $\pi \in \P_{1,2}^{B}(n)$ there exists a unique {reflection block} $\bar B\in\pi$.
This leads to one more definition.
We call $\BL= ((\bar b, \bar a),( a,b))$ a \emph{B-pair} if $\bar b <b$, $|{a}| <|b|$ and $(\bar b,\bar a ), (a,b)\in \pi$.
The B-pair $((\bar b, \bar a),( a,b))$ is called \emph{positive} if $b\times a >0 $; otherwise
it is called \emph{negative}; see Figure~\ref{fig:coneced}.
The set of such B-pairs is denoted by $\Pair_B(\pi)$.
Let us notice that from Definition~\ref{def:partycji} it follows that
the element $s$ is a singleton of $\pi \in \P_{1,2}^{B}(n)$ if and only if $\bar s$ is also a singleton of $\pi$, thus we can define the subset of B-singletons $((\bar a),(a))$ as
\[
\Sing_B(\pi):=\{((\bar a), (a)) \mid \bar a,a \in \Sing(\pi), \,\bar a < a\}, \qquad \pi \in \P_{1,2}^{B}(n).
\]
Let
\[
\P_{2}^{B}(n):=\big\{\pi \in \P_{1,2}^{B}(n)\mid \Sing_B(\pi)=\varnothing\big\}
\]
and
\[
\PA_2(n):=\big\{\pi \in \P_{2}^{B}(n)\mid (\bar b,\bar a ),\, (a,b)\in \pi \implies ((\bar b,\bar a ), (a,b))\text{ is positive}\big\},
\]
i.e., it is a subset of $\P_{2}^{B}(n)$, with only positive B-pairs.

\begin{Remark}\quad
\begin{enumerate}
\item [$1.$] We note that a B-pair is not a block of $\pi$, but a pair of reflection pairs.
	
\item [$2.$] We note that $\#\PA_2(n)=(n-1)!!$ for even $n$, which is the same as the number of the classical pair partitions on $n$ points.
	In Figure~\ref{figexample1}, they are depicted on the top row.
\end{enumerate}	
\end{Remark}

\begin{figure}[h]\centering
\MatchingMeandersabc{4}{-4/1, -3/-2,2/3,-1/4 }{}
\caption{The example of $\pi \in \P_{2}^{B}(4)$ with a positive B-pair $((\bar 3,\bar 2 ), (2,3))$ and a negative B-pair $((\bar 4,1), (\bar 1,4))$. }
	\label{fig:coneced}
\end{figure}

We introduce some partition statistics.
For two pairs $V$ and $ W$, we introduce the relation $\text{cr}$ as follows:
\begin{gather*}
V\stackrel{\text{cr}}{\sim}W \iff
V=(i,j),\qquad W=(k,l)\quad \text{such that}\quad i<k<j<l,
\end{gather*}
For a set partition $\pi\in \P_{1,2}^{B}(n)$, let $\Cr(\pi)$ be the number of crossings of B-pairs, i.e.,
\begin{align*}
\Cr(\pi)={}&\#\big\{\big(\big(\overline{V}, V \big),\big(\overline{ W}, W\big)\big)\in \Pair_B(\pi)\times \Pair_B(\pi) \mid \text{$V\stackrel{\text{cr}}{\sim}W $}\big\}
\\
&+\#\big\{\big(\big(\overline{V}, V\big),\big(\overline{ W}, W\big)\big)\in \Pair_B(\pi)\times \Pair_B(\pi) \mid \text{$\overline{V}\stackrel{\text{cr}}{\sim}W $}\big\}.
\end{align*}
For two blocks $V$, $W$ of a set partition, we say that $W$ \emph{covers} $V$ if there are $i,j \in W$ such that $i <k <j$ for any $k\in V$ and then we write $V\stackrel{\text{cs}}{\sim}W $. For $\pi\in\P_{1,2}^{B}(n)$, let $\InS(\pi)$
be the number of pairs of a B-singleton and a covering B-pair:
\begin{align*}
\begin{split}
\InS(\pi)&=\#\big\{\big(\big(\overline{ V}, V \big),\big(\overline{ W}, W\big)
\big) \in \Sing_B(\pi) \times \Pair_B(\pi) \mid V\stackrel{\text{cs}}{\sim}W
\big\}
\\
 & +\#\big\{\big(\big(\overline{V}, V \big),\big(\overline{W}, W\big)\big) \in \Sing_B(\pi) \times \Pair_B(\pi) \mid \overline{V}\stackrel{\text{cs}}{\sim}W\big\}.
 \end{split}
\end{align*}
We say that the B-pair $(\overline{ W} , W)$ cover the B-singleton $\big(\overline{V}, V\big)$ if $V\stackrel{\text{cs}}{\sim}{W} $ or $\overline{V}\stackrel{\text{cs}}{\sim}{ W} $ (equivalently, $\big(\overline{V}, V\big)$ is the inner singleton of $\big(\overline{W} , W\big)$).

Let $\NB(\pi)$ be the number of negative B-pairs of $\Pair_B(\pi)$, where $\pi \in \P_{1,2}^{B}(n)$.
A set partition $\pi \in \P_{1,2}^{B}(n)$ is \emph{non-crossing} if $\Cr(\pi)=0$. The set of non-crossing pair partitions of $[\pm n]$ is denoted by $\NC^B_2(n)$, and by $\NC^A_2(n)$ we denote the subset of $\NC^B_2(n)$ where all B-pairs are positive.

\begin{Remark} \label{jakliczyc}
While reading this, the first impression seems to be that the procedure of counting the crossings is complicated. This is not true since it can be read from the figure as follows. First, we draw a vertical line in the center as in Figure~\ref{fig:examplejakliczyccrosingi}.
Then we count only the crossings on the left of this vertical line.
We count the number of
negative B-pairs as the number of B-pairs crossed by a vertical line ($\scriptstyle{\blacklozenge}$ points on Figure~\ref{fig:examplejakliczyccrosingi}).
Similarly we count the pairs of the form
\[
(\text{B-singleton, covering B-pair})
\]
as the number of such a pairs with at least one leg on the left of this vertical line.

\begin{figure}[htp]
\centering
\MatchingMeanders{10}{-1/3, -3/1 ,-6/5,-5/6,7/10,-10/-7,-2/9,-9/2}
\caption{The example of statistic of partition $\pi \in \P_{1,2}^{B}(10)$, i.e., $\Cr(\pi)=4$, $\NB(\pi)=3$, $\InS(\pi)=5$.}
\label{fig:examplejakliczyccrosingi}
\end{figure}
\end{Remark}

Now we prove the following theorem, which shows the relationship between the set of partitions of type B
(corresponding statistic) and a joint action of Gaussian operators on a vacuum vector.

\begin{Theorem}\label{lem101}
For any $i\in\{1,\ldots,2n\}$ and $x_{\bar i}\otimes x_i \in \HH_\R$, we have
\begin{equation}\label{formula101}
	\state(\G(x_{\overline{ 2n}}\otimes x_{ 2n})\cdots \G(x_{\bar 1}\otimes x_{ 1}))= \sum_{\pi\in\PB_{2}(2n)} \s^{\NB(\pi)}\q^{\Cr(\pi)} \prod_{\substack{(i,j) \in \Pair(\pi)} }\langle x_i, x_j\rangle.
\end{equation}
\end{Theorem}


\begin{proof}
Given
$\epsilon=(\epsilon(1), \dots, \epsilon(n))\in\{1,\ast\}^{n}$, let
$\P_{1,2;\epsilon}^{B}(n)$ be the set of partitions $\pi\in
\P_{1,2}^{B}(n)$ such that
\begin{itemize}\itemsep=0pt
\item if $((\bar b ,\bar a) , (a,b))$ is a B-pair of
$\Pair_B(\pi)$, then $\epsilon(|a|)= \ast$, $\epsilon(|b|)=1$,
\item if $\{c\}$ is a singleton in $\pi$, then $\epsilon( |c|)=\ast$.
\end{itemize}
We will prove first that
\begin{gather}
\B^{\epsilon(n)}(x_{\overline{n}}\otimes x_{n})\cdots \B^{\epsilon(1)}(x_{\overline{1}}\otimes x_1)\Omega \otimes \Omega \nonumber
\\
\qquad{} = \sum_{\pi\in\PB_{1,2;\epsilon}(n)} \s^{\NB(\pi)}q^{\Cr(\pi)+ \InS(\pi)}
\prod_{\substack{(i,j) \in \Pair(\pi)} }\langle x_i, x_j\rangle
\bigotimes_{i\in \Sing(\pi)} x_i\label{formula10112}
\end{gather}
holds, where
\[
\Sing(\pi) = \{(\bar v_m),\dots,(\bar v_1), (v_1),\dots,(v_m)\} \subset\bar{\N} \cup \N, \qquad
\bar v_m<\cdots < \bar v_1< v_1<\cdots < v_m,
\]
and
$\bigotimes_{i\in \Sing(\pi)} x_i$ denotes the tensor product $x_{\bar v_m}\otimes \cdots \otimes x_{\bar v_1} \otimes x_{v_1}\otimes \cdots \otimes x_{v_m}$.
If
\[
\#\{i\in[j] \mid \epsilon(i)=1\}>\#\{i\in[j] \mid \epsilon(i)=\ast\}
\]
for some $j\in\{1,\dots,n\}$, then we understand the sum over the empty set is 0 since $\PB_{1,2;\epsilon}(n)=\varnothing$ in this case.
The proof of \eqref{formula10112} is given by induction.

When $n=1$, then $\B(x_{\bar 1}\otimes x_1)\Omega\otimes \Omega=0$ and $\B^\ast(x_{\bar 1}\otimes x_1)\Omega\otimes \Omega=x_{\bar 1}\otimes x_1$ and hence the formula \eqref{formula10112} is true.

Suppose that the \eqref{formula10112} is true for $n = k$.
We will show that the action of $\B^{\epsilon(k+1)}(x_{\,\overline{k+1}}{}\otimes x_{k+1})$ for $\epsilon(k+1)=\ast$ or $\epsilon(k+1)=1$ corresponds to the inductive pictorial description of set partitions of type B.

\smallskip\noindent
\emph{Case} 1. If $\epsilon(k+1)=\ast$, then the operator $\B^\ast(x_{\,\overline{k+1}}\otimes x_{k+1})$ acts on the tensor product, putting~$x_{k+1}$ on the right and $x_{\,\overline{k+1}}$ on the left. This operation pictorially corresponds to adding the singleton $(\overline{k+1}
)$ and $(k+1)$ to $\pi\in\PB_{1,2;\epsilon}(k)$, to yield the new type-B partition $\tilde{\pi}\in\PB_{1,2;\epsilon}(k+1)$ such that $((\overline{k+1}), (k+1))\in \Sing_B(\tilde{\pi})$. This map $\pi\mapsto \tilde{\pi} $ does not change the numbers $\NB$, $\Cr$ or $\InS$, which is compatible with the fact that the action of $\B^\ast(x_{\,\overline{k+1}}\otimes x_{k+1})$ does not change the coefficient.
Hence the formula~\eqref{formula10112} holds when $n=k+1$ and $\epsilon(k+1)=\ast$.

\smallskip\noindent
\emph{Case} 2. If $\epsilon(k+1)=1$, then we have two subcases, depending on $p_\q$ and $\ell_{\q}$.

If the positive part $p_\q $ (subcase (a)) (resp. $\s \ell_{\q}$ -- subcase (b)) acts on the tensor product on the right hand side of \eqref{formula10112} (for fixed $\pi$), then new $m$ terms appear by using \eqref{rq}. We obtain the equations
\begin{subequations}
\begin{align}
	\begin{split}
		p_\q(x_{\,\overline{k+1}}\otimes x_{k+1}) & x_{\bar v_m}\otimes \cdots \otimes x_{v_m}=
		\sum_{i=1}^m \q^{m-i}\langle x_{\,\overline{k+1}},x_{\bar v_i}\rangle \langle x_{v_i} ,x_{k+1}\rangle\, \\ &\times x_{\bar v_m }\otimes \cdots \otimes \check{x}_{\bar v_i} \otimes \cdots \otimes x_{\bar v_1} \otimes x_{v_1}\otimes \cdots \otimes \check{x}_{v_i} \otimes \cdots \otimes x_m,
	\end{split}\label{eq:equatinpomocdowodgolwny1}
	\\
	\begin{split}
		\alpha\ell_\q(x_{\,\overline{k+1}}\otimes x_{k+1}) & x_{\bar v_m}\otimes \cdots \otimes x_{v_m}=
		\alpha\q^{m-1}\sum_{i=1}^m \q^{i-1}
		\langle x_{\,\overline{k+1}},x_{ v_i}\rangle \langle x_{\bar v_i} ,x_{k+1}\rangle\, \\& \times x_{\bar v_m }\otimes \cdots \otimes \check{x}_{\bar v_i} \otimes \cdots \otimes x_{\bar v_1} \otimes x_{v_1}\otimes \cdots \otimes \check{x}_{v_i} \otimes \cdots \otimes x_m.
	\end{split}\label{eq:equatinpomocdowodgolwny2}
\end{align}
\end{subequations}
Now we will focus on the $i$-th summand of the above equations.

We fix $\pi\in \PB_{1,2;\epsilon}(k)$ and suppose that $\pi$ has singletons
\begin{align*}
\bar v_m<\cdots <\bar v_{ i} < \cdots <\bar v_1 < v_1<\cdots <v_i <\cdots <v_m,\text{ where } i\in [m].
\end{align*}

We assume that $\Pair_B(\pi)$ contains the B-pairs $\big(\overline{ W}_1 , W_1\big),\dots, \big(\overline{ W}_u , W_u\big)$ which cover the B-singleton $((\bar v_i),( v_i))$.
Case~2(a)~-- see Figure~\ref{figtracepositivea}.
In the $i$-th term of \eqref{eq:equatinpomocdowodgolwny1} the inner product $\langle x_{\,\overline{k+1}},x_{\bar v_i}\rangle \langle x_{v_i} ,x_{k+1}\rangle$ appears with coefficient $\q^{m-i}$. Pictorially this corresponds to getting a set partition $\tilde{\pi} \in \PB_{1,2;\epsilon}(k+1)$ by adding the positive B-pair $\big((\overline{k+1}, \overline{v}_i) , (v_i, k+1)\big)$ to $\Pair_B(\tilde \pi)$. This new B-pair crosses the B-pairs $\big(\overline{W}_j , W_j\big)$, $j\in\{1,\dots,u\}$ and so increases the crossing number by~$u$ but decreases the number $\InS(\pi)$ by $u$ because originally $((\bar v_i),( v_i))$ was the inner singleton of B-pairs $\big(\overline{W}_j , W_j\big)$, $j\in\{1,\dots,u\}$.
Now the new inner B-singletons $((\bar v_{i+1}), (v_{i+1})), \dots,(( \bar v_m), (v_m))$ of B-pair $((\overline{k+1}, \overline{v}_i) , (v_i, k+1))$
appear.
Altogether we have
\[
\Cr(\tilde{\pi})=\Cr(\pi)+u, \qquad
\InS(\tilde{\pi})=\InS(\pi)+m-i-u\qquad \text{and}\qquad
\NB(\tilde{\pi})=\NB(\pi).
\]
So the exponent of $\q$ increases by $m-i$. This factor $\q^{m-i}$ is exactly the factor appearing in equation~\eqref{eq:equatinpomocdowodgolwny1}.

\begin{figure}[h]
\centering
\MatchingMeandersPositiveTracea{9}{-9/-8,8/9,-6/2,-2/6 }{}
\MatchingMeandersPositiveTraceb{10}{-9/-8,8/9,-6/2,-2/6 ,-10/-4,4/10 }{}
\caption{The visualization of the action of $p_\q(x_{\,\overline{k+1}}\otimes x_{k+1})$.}
\label{figtracepositivea}
\end{figure}

Case 2(b) -- see Figure~\ref{figtracepositive}.
In the $i$-th term of \eqref{eq:equatinpomocdowodgolwny2} the inner product $\langle x_{\,\overline{k+1}},x_{ v_i}\rangle \langle x_{\bar v_i} ,x_{k+1}\rangle$ appears with the coefficient $\s \q^{m+ i-2}$.
Graphically this
corresponds to getting a set partition $\tilde{\pi} \in \PB_{1,2;\epsilon}(k+1)$ by adding $\overline{k + 1}$ and $k + 1$ to $\pi$ and creating the new negative B-pair $\big((\overline{k+1}, v_{i}) , (\overline v_ i, k+1)\big)\in \Pair_B(\tilde \pi)$.
Similarly to Case 2(a), we count the change of numbers and get
\[
\Cr(\tilde{\pi})=\Cr(\pi)+u,\quad
\InS(\tilde{\pi})=\InS(\pi)+m-1+i-1-u\quad\text{and}\quad \NB(\tilde{\pi})=\NB(\pi)+1.
\]
Altogether, when moving from $\pi$ to $\tilde{\pi}$, the exponent of $\s$ increases by $1$ and the exponent of $\q$ increases by ${m+ i-2}$, which coincides with the coefficient appearing in the action of $\s \ell_{\q}(x_{\,\overline{k+1}}\otimes x_{k+1})$, creating the inner product $\langle x_{\,\overline{k+1}},x_{ v_i}\rangle \langle x_{\bar v_i}, x_{k+1}\rangle$.

\begin{figure}[h]
%\begin{center}
%\QQ{$(b)$}
%\end{center}%
\centering
\MatchingMeandersPositiveTracea{9}{-9/-8,8/9,-6/2,-2/6 }{} % $\xrightarrow{\text{trace action }}$
\MatchingMeandersPositiveTraceb{10}{-9/-8,8/9,-6/2,-2/6 ,-10/4,-4/10 }{}
\caption{The visualization of the action of $\s \ell_{\q}(x_{\,\overline{k+1}}\otimes x_{k+1})$.}
\label{figtracepositive}
\end{figure}


Note that as $\pi$ runs over $\PB_{1,2;(\epsilon(1),\dots,\epsilon(k))}(k)$, every set partition $\tilde{\pi}\in \PB_{1,2;(\epsilon(1),\dots,\epsilon(k),1)}(k+1)$ appears exactly once, either in Case~2(a) or in Case~2(b). Therefore, in Case~2, the pictorial inductive step and the actual action of $\B(x_{\,\overline{k+1}}\otimes x_{k+1})$ both create the same terms with the same coefficients, and hence the formula \eqref{formula10112} is true when $n=k+1$ and $\epsilon(k+1)=1$. Case~1 and Case~2 show by induction that the formula \eqref{formula10112} holds for all $n \in \N$.

For a given set of B-pairs $\{((\bar b_1, \bar a_1),(a_1,b_1)), \dots, ((\bar b_{n/2}, \bar a_{n/2}),(a_{n/2},b_{n/2}))\}\in \Pair_B( \pi) $,
where $\pi \in \PB_{2}(n)$, we denote the set of \emph{left and right legs} of $\pi$ by $l_\pi :=\{ a_1,\bar a_1,\dots, a_{n/2}, \bar a_{n/2} \}$, and $r_\pi :=\{ b_1,\bar b_1,\dots, b_{n/2}, \bar b_{n/2} \}$. Formula \eqref{formula101} follows from \eqref{formula10112} by taking the sum over all $\epsilon$ such that $\Sing(\pi)=\varnothing$. In this case, we understand that $\bigotimes_{i\in \Sing(\pi)} x_i=\Omega\otimes \Omega $ and
\begin{align*}
\bigsqcup_{\e\in\{1,\ast\}^{n}} \PB_{2;\e}(n)&=\bigsqcup_{\e\in\{1,\ast\}^{n}} \big\{\pi \in \P_{2}^{B}(n)\mid (\bar b,\bar a ), (a,b)\in \pi \implies \epsilon(|a|)= \ast,\, \epsilon(|b|)=1\big\}
\\
& =\bigsqcup_{\substack{ L\subset [\pm n] \\ \# L =n
}} \big\{\pi \in \P_{2}^{B}(n)\mid
l_\pi=L \text{ and }r_\pi = [\pm n]\setminus L\big\}
\\
&=\PB_{2}(n).
\end{align*}

 Finally, applying the state action, we get
\begin{align*}
&\state(\G(x_{\overline{ 2n}}\otimes x_{ 2n})\cdots \G(x_{\bar 1}\otimes x_{ 1}))
\\
&\qquad =\sum_{\e\in\{1,\ast\}^{2n}} \state\big( \B^{\epsilon(2n)}(x_{\overline{2n}}\otimes x_{2n})\cdots \B^{\epsilon(1)}(x_{\overline{1}}\otimes x_1)\big)
\\
&\qquad =\sum_{\e\in\{1,\ast\}^{2n}} \sum_{\substack{\pi\in\PB_{1,2;\epsilon}(2n)\\ \Sing(\pi)=\varnothing }} \s^{\NB(\pi)}q^{\Cr(\pi)+ \InS(\pi)}
\prod_{\substack{(i,j) \in \Pair(\pi)} }\langle x_i, x_j\rangle
\\
&\qquad =\sum_{\e\in\{1,\ast\}^{2n}} \sum_{\substack{\pi\in\PB_{2;\epsilon}(2n) }} \s^{\NB(\pi)}q^{\Cr(\pi)}
\prod_{\substack{(i,j) \in \Pair(\pi)} }\langle x_i, x_j\rangle
\\
&\qquad = \sum_{\pi\in\PB_{2}(2n)} \s^{\NB(\pi)}\q^{\Cr(\pi)} \prod_{\substack{(i,j) \in \Pair(\pi)} }\langle x_i, x_j\rangle.
\tag*{\qed}
\end{align*}
\renewcommand{\qed}{}
\end{proof}
\begin{thebibliography}{99}
\bibitem[9]{BDEG2021}
Bo\.zejko M., Do{\l}\c{e}ga M., Ejsmont W., Gal S.R., Reflection length with
 two parameters in the asymptotic representation theory of type~{B}/{C} and
 applications, \href{https://doi.org/10.1016/j.jfa.2022.109797}{\textit{J.~Funct. Anal.}} \textbf{284} (2023), 109797, 47~pages,
 \href{https://arxiv.org/abs/2104.14530}{arXiv:2104.14530}.

\bibitem[10]{BozejkoEjsmontHasebe2015}
Bo\.zejko M., Ejsmont W., Hasebe T., Fock space associated to {C}oxeter groups
 of type~{B}, \href{https://doi.org/10.1016/j.jfa.2015.06.026}{\textit{J.~Funct. Anal.}} \textbf{269} (2015), 1769--1795,
 \href{https://arxiv.org/abs/1411.7997}{arXiv:1411.7997}.

\bibitem[13]{BozejkoSpeicher1994}
Bo\.zejko M., Speicher R., Completely positive maps on {C}oxeter groups,
 deformed commutation relations, and operator spaces, \href{https://doi.org/10.1007/BF01450478}{\textit{Math. Ann.}}
 \textbf{300} (1994), 97--120, \href{https://arxiv.org/abs/funct-an/9408002}{arXiv:funct-an/9408002}.

\bibitem[14]{BozejkoSzwarc2003}
Bo\.zejko M., Szwarc R., Algebraic length and {P}oincar\'e series on reflection
 groups with applications to representations theory, in Asymptotic
 Combinatorics with Applications to Mathematical Physics ({S}t.~{P}etersburg,
 2001), \textit{Lecture Notes in Math.}, Vol.~1815, \href{https://doi.org/10.1007/3-540-44890-X_9}{Springer}, Berlin, 2003,
 201--221.

\bibitem[20]{Ejsmont2020}
Ejsmont W., Fock space associated with quadrabasic {H}ermite orthogonal
 polynomials, \href{https://doi.org/10.4064/sm210916-19-8}{\textit{Studia Math.}} \textbf{270} (2023), 17--56,
 \href{https://arxiv.org/abs/2004.08538}{arXiv:2004.08538}.

\end{thebibliography}
\end{document}
